\section{The matching condition and further questions}\label{sec:comparison}

\subsection{An edgewise switching example}
The following calculation explains the distinction between the uniqueness
proved here and the edgewise conditions printed in
\cite[Definition~7.8]{HKKP26}. It concerns those displayed conditions;
it is not asserted to define a spectral-network object in the sense of
\cite{HKS24}.

Fix $p>0$, put $r=(-p,0,p)$ and $\tau=3$, and let
$I=(-2,-1)$, $J=(1,2)$. Let $\sigma$ exchange indices $1,3$
and fix $2$. For $L=I$ or $J$ and $t\le\tau$, define
\begin{equation}\label{eq:switchweights}
 w_i^L(t)=\int_{-\infty}^t
 \left[\boldsymbol1_{s\notin L}\lambda(r_i,s)
       +\boldsymbol1_{s\in L}\lambda(r_{\sigma(i)},s)\right]\,\dd s.
\end{equation}
Use these weights at both vertices in the standard basis.
The identity arrow is an isometry and the norms tend to the
same standard incoming norm.

The function $\lambda(p,s)$ is even in $s$. Its integrals over
$I$ and $J$ agree. Since the first and third positions are
opposites and the second is zero, the two sets of weights in
\eqref{eq:switchweights} coincide for $2\le t\le3$.
Write $k_i$ for their common value at $t=\tau$. Continue both
paths by
\[
 w_{vi}(t)=k_i+
   \int_\tau^t\lambda(r_i+\epsilon_va(s-\tau),s)\,\dd s,
 \qquad t\ge\tau.
\]
The coordinate gaps are positive for $t>\tau$, and the
outgoing derivatives agree with the incoming derivatives
at $\tau$. Both normalized outgoing limits exist, by the
same integrable difference estimate as in
\eqref{eq:scalarboundarymap}. Since the paths coincide on
this whole outgoing tail, those full limits are equal.

The two paths are nevertheless different. On $-2<t<-1$,
\[
 w_1^I(t)-w_1^J(t)
       =2\int_{-2}^t\lambda(p,s)\,\dd s>0.
\]
Evaluation on $e_1$ detects their difference as actual norms.
At every regular incoming time the unordered positions of
either path are $\{-p,0,p\}$, with stable $B$ labels.
At each switch the coordinate subspaces carrying those
positions are exchanged. After $\tau$ their supports are
the same three Ys and their crossings, translated in time
from \eqref{eq:diamondvertices}. This finite support is
connected and balanced. Subdivide at the switches and
crossings; all open edges have stable labels and the required
straight directions. At a switch both adjacent abstract
labels are $B$, and numerical charge conservation holds.

At a switch, however, the one-sided coordinate velocity
vectors in the fixed apartment differ. In the $u$ coordinate,
if they are $v^-,v^+$ at a switch $u_0$, then the Euclidean
apartment formula gives
\[
 d_2\bigl(y_v(u_0),\midpt(y_v(u_0-h),y_v(u_0+h))\bigr)
       =\tfrac h2\norm{v^+-v^-}_2+O(h^2).
\]
The coefficient is nonzero. Thus both paths fail
\eqref{eq:matching}. Equality of support, phases, abstract
edge labels and full ends has not identified the embedded
velocity subspaces across the exceptional times.

The printed definition describes an edge label as a graded
object of a velocity filtration. If that phrase is intended
to include additional transport of its embedding, that
transport is precisely what must be specified to exclude
this example. We have made one sufficient intrinsic condition
explicit and proved the correspondence in the resulting
finite regular class. The calculation above does not disprove
an interpretation incorporating stronger categorical gluing.

\subsection{Comparison with categorical networks}
Proposition~\ref{prop:vertexfiltrations} goes from matched norm
paths to local residue objects, two increasing-phase filtrations
and their edge quotients. The reverse implication is a separate
question: which globally framed network data determine opposite
first-order norm germs? In particular, the local existence of
two filtrations of an abstract representation does not by itself
supply the embedded opposition used in
\eqref{eq:vertexdecomposition}.

A comparison with the global Fukaya-category construction would
also have to identify the relevant Novikov-field object and its
gluing data. The uniqueness theorem of
\cite[Theorem~4.15]{HKS24}, in its constant-coefficient,
product-type setting, gives support equality and S-equivalence
of reductions for networks representing the same object after
Novikov extension. It does not supply this comparison or assert
equality of full norm paths. The present variational uniqueness
uses the specified boundary triple and the midpoint condition
directly.

\subsection{Dependence on the boundary data}
Theorem~\ref{thm:main} includes every full split boundary triple,
and hence includes degenerate inputs. General continuity of the
assignment remains open here. A natural topology on triples is
the maximum of the three unnormalized $d_2$ distances. For paths,
one may use uniform vertexwise $d_2$ distance after applying the
same prescribed phase scalars on each end. The estimate in
Proposition~\ref{prop:collapsecontinuity} establishes this
continuity for one explicit family through a bounded-face
collapse. A general result would have to control comparisons
between variational problems with different full reference
norms, rather than only competitors in one fixed-end problem.

Two other extensions are distinct from that continuity problem.
First, the theorem fixes the symmetric $A_2$ charges and radial
area form; its finite-support proof uses the three resulting
monotone directions. Second, uniqueness has been proved among
finite regular paths with matching. Extending it to a larger
intrinsically defined class would require a regularity theorem
for that class or a suitable direct comparison, even though the
minimum itself already has a finite algebraic presentation.
